% ============ 02. Introducción y alcance ============
\section{Introducción, alcance y método de la revisión}

\subsection{Propósito}

Este documento es un \textbf{estado del arte} orientado a la mejora del
repositorio \texttt{Observatorio\_Ministerio\_de\_Ciencias\_Grupo8} (observatorio
crítico del sistema de reconocimiento de investigadores de MinCiencias,
Colombia). Su objetivo es doble: (i) situar el proyecto en la literatura
científica relevante — qué se sabe, qué métodos se recomiendan y qué han hecho
otros autores en problemas equivalentes — y (ii) traducir esa literatura en
\textbf{acciones concretas de mejora} para este repositorio, respetando su
alcance (6 convocatorias 2013--2021, 77\,237 registros declarados, 30\,086
investigadores únicos, cruce con 3\,166\,629 productos de grupos) y sus
limitaciones documentadas en la auditoría (consolidado versionado incompleto,
déficit de inferencia formal, reproducción end-to-end pendiente).

\subsection{Criterio de búsqueda y fuentes}

La revisión \textbf{no se limita por idioma}: incorpora literatura en inglés
(núcleo de la bibliometría y la ciencia de la ciencia), español y portugués
(estudios sobre los sistemas de información curricular de Brasil y Colombia,
evaluación nacional, ciencia abierta) y materiales grises relevantes
(declaraciones y acuerdos de política: DORA, Leiden, CoARA). La ventana
prioritaria es 2018--2026, con los clásicos fundacionales anteriores cuando
siguen siendo la referencia canónica (Lotka 1926; Hirsch 2005; Wang y Strong
1996; Kimball 2013; Huang 1998).

\subsection{Cómo leer este documento}

Cada sección temática sigue el mismo patrón: \textbf{(a) qué dice la
literatura} (autores y hallazgos), \textbf{(b) cómo se relaciona con este
repositorio} (estado actual verificado en la auditoría) y \textbf{(c) mejoras
recomendadas}, con referencias. La sección 12 consolida todo en una matriz de
mejora accionable. Las referencias canónicas aparecen en la bibliografía
(sección 13.1); los recursos descubiertos por búsqueda en línea, citados con
su URL (sección 13.2), sin fabricar metadatos de autoría no verificados.

\begin{tcolorbox}[nota]
\textbf{Glosario rápido} (lenguaje llano): \emph{bibliometría}, estudio
cuantitativo de publicaciones y citas; \emph{ciencia de la ciencia (SciSci)},
disciplina que estudia el sistema científico con datos masivos; \emph{métricas
responsables}, uso de indicadores cuantitativos solo como apoyo al juicio;
\emph{atrición}, abandono del sistema científico; \emph{provenance/linaje},
trazabilidad de cómo se produjo cada dato.
\end{tcolorbox}
